% =====================================================================
% Kapitel 3: System und Pipeline
% =====================================================================

\chapter{\agora{}: System und Pipeline}\label{ch:system}

\section{Architekturüberblick}\label{sec:architektur}

\agora{} ist eine lokal oder kontrolliert hybrid betreibbare Multi-Agent-Plattform für simulierte DACH-Zielgruppen-, Stakeholder- und Marktreaktionen. Die Plattform ist als Flask/Python-3.14-Backend mit Neo4j-Graph, Redis-Cache, OASIS- bzw.\ CAMEL-Simulations-Subprozess und Vue-3-Frontend strukturiert. Die Code-Basis umfasst zum Auditzeitpunkt etwa 3500~LOC im Backend-Service-Layer, 5500~LOC im OASIS-Subprozess und etwa 3900~LOC Graph-bezogenen Code. Der Reifegrad ist \emph{0.9.3 Stability Beta}; das Ziel ist nach \code{ROADMAP.md} eine stabile Single-User-Version \code{1.0.0}\footnote{\raggedright Status-Informationen aus~\code{AGENTS.md} und \code{ROADMAP.md} im Repository.}

Für die vorliegende Arbeit sind zwei Schichten zentral: die \textbf{Evidence-Pipeline} (vom Seed-Dokument zum Claim) und die \textbf{Simulations-Pipeline} (Persona-Population, OASIS-Aktionen, Echo-Index). Beide sind über die \code{EvidenceSourceKind}-Taxonomie~\cite{adr-0011} miteinander verbunden.

\section{Die Evidence-Chain}\label{sec:evidencechain}

Abbildung~\ref{fig:evidence-chain} rekonstruiert die Evidence-Chain vom Seed-Dokument bis zum ausgelieferten Claim in 17 Stufen. Die Rekonstruktion beruht auf statischer Code-Analyse von~\code{backend/app/services/} und~\code{backend/app/contracts/}, validiert durch die Code-Review-Graph-Werkzeuge. Sieben Stufen (\textsc{✗}) markieren code-verifizierte Provenance-Verlustpunkte.

\begin{figure}[ht]
  \centering
  % =====================================================================
% figures/evidence-chain.tex
% TikZ-Rekonstruktion der Evidence-Chain (Audit §1, vereinfacht)
% =====================================================================

\begin{tikzpicture}[
  font=\small,
  stage/.style={draw, rounded corners=2pt, minimum width=44mm, minimum height=9mm,
                align=center, fill=blue!5},
  losspoint/.style={draw, rounded corners=2pt, minimum width=44mm, minimum height=9mm,
                    align=center, fill=red!8, drop shadow},
  arrow/.style={-{Stealth[length=2.5mm]}, thick, draw=blue!60!black},
  lossyearrow/.style={-{Stealth[length=2.5mm]}, thick, draw=red!70!black, dashed},
  label/.style={font=\footnotesize, fill=white, inner sep=1pt},
  losslabel/.style={label, draw=red!70!black, rounded corners=1pt, inner sep=1.5pt}
]

% Hauptkette (Stufen 1-17)
\node[stage] (s1)  {[1] Seed-Dok.\\(Upload: PDF/MD/TXT)};
\node[stage, below=8mm of s1]  (s2)  {[2] Chunks\\\texttt{split\_text}};
\node[losspoint, below=8mm of s2]  (s3)  {\textbf{[3] KG}\\Neo4j \texttt{add\_text}};
\node[stage, below=8mm of s3]  (s4)  {[4] Retrieval\\\texttt{search\_graph}};
\node[losspoint, below=8mm of s4]  (s5)  {\textbf{[5] Report Tools}\\\texttt{\_record\_tool\_evidence}};
\node[stage, below=8mm of s5]  (s6)  {[6] Persona/Sim\\OASIS/CAMEL};
\node[stage, below=8mm of s6]  (s7)  {[7] Interview\\\texttt{AgentInterview}};
\node[stage, below=8mm of s7]  (s8)  {[8] Evid.-Norm.\\\texttt{\_TYPE\_TO\_SOURCE\_KIND}};
\node[stage, below=8mm of s8]  (s9)  {[9] Evid.-Ident.\\\texttt{build\_evidence\_id}};
\node[stage, below=8mm of s9]  (s10)  {[10] Evid.-Index\\\texttt{evidence-map.json}};
\node[stage, below=8mm of s10] (s11) {[11] Claim-Bind.\\\texttt{evidence\_binder}};
\node[stage, below=8mm of s11] (s12) {[12] Gating\\(ADR-0002)};
\node[stage, below=8mm of s12] (s13) {[13] Confidence\\\texttt{compute\_conf.}};
\node[stage, below=8mm of s13] (s14) {[14] Hypothesen\\\texttt{dedup\_and\_cap}};
\node[stage, below=8mm of s14] (s15) {[15] Red-Team\\(LLM-Judge)};
\node[losspoint, below=8mm of s15] (s16) {\textbf{[16] Renderer}\\\texttt{markdown\_renderer}};
\node[stage, below=8mm of s16] (s17) {[17] Export\\(JSON/ZIP/MD/CSV)};

% Pfeile
\draw[arrow] (s1) -- (s2);
\draw[lossyearrow] (s2) -- (s3) node[midway, losslabel, xshift=3mm] {\lossmark{} 1};
\draw[lossyearrow] (s3) -- (s4) node[midway, losslabel, xshift=3mm] {\lossmark{} 2,3,4};
\draw[lossyearrow] (s4) -- (s5) node[midway, losslabel, xshift=3mm] {\lossmark{} 5};
\draw[lossyearrow] (s5) -- (s6) node[midway, losslabel, xshift=3mm] {\lossmark{} 6,7};
\draw[arrow] (s6) -- (s7);
\draw[arrow] (s7) -- (s8);
\draw[arrow] (s8) -- (s9);
\draw[arrow] (s9) -- (s10);
\draw[arrow] (s10) -- (s11);
\draw[arrow] (s11) -- (s12);
\draw[arrow] (s12) -- (s13);
\draw[arrow] (s13) -- (s14);
\draw[arrow] (s14) -- (s15);
\draw[lossyearrow] (s15) -- (s16);
\draw[arrow] (s16) -- (s17);

% Legende (inline, ohne tikz-matrix)
\node[stage, minimum width=22mm, minimum height=5mm, below=10mm of s17, xshift=-30mm] (leg1) {};
\node[losspoint, minimum width=22mm, minimum height=5mm, right=2mm of leg1] (leg2) {};
\node[label, right=2mm of leg2] {\lossmark{} = Provenance-Verlustpunkt};

\end{tikzpicture}

  \caption{Rekonstruierte Evidence-Chain vom Seed-Dokument bis zum Claim. Die sieben mit \textsc{✗} markierten Stufen sind code-verifizierte Provenance-Verlustpunkte.}
  \label{fig:evidence-chain}
\end{figure}

\section{Die \code{EvidenceSourceKind}-Taxonomie}\label{sec:sourcekind}

Der Vertrag~\code{backend/app/contracts/report\_contract.py} definiert sechs Werte für den Enum~\code{EvidenceSourceKind}:

\begin{itemize}
  \item \code{seed\_corpus} -- echte externe Quelle aus dem hochgeladenen Korpus. Aktuell nur als Zielzustand in~\cite{adr-0013} spezifiziert; im Referenzlauf~\cite{audit-r2} nicht erreicht.
  \item \code{agent\_quote} -- Persona-Interview-Zitat. Provenance über~\code{persona\_id}, \code{stakeholder\_group}, \code{quote}.
  \item \code{agent\_action} -- sozialer Akt aus der Simulation (Post, Comment, Reshare). Provenance über \code{sim\_action\_log} (Plattform, Runde, Agent, Aktion, Zeitstempel).
  \item \code{graph\_relation} -- Fakt aus dem Neo4j-Graph. Aktuell nur syntaktisch referenzierbar; \emph{keine} Dokumentidentität.
  \item \code{web\_source} -- URL aus externer Recherche. URL im Anker; Fetch ist echtzeitabhängig und kann verfallen.
  \item \code{inferred} -- Standard-Fallback, wenn die Herkunft unbekannt ist. \code{Confidence}-Gewicht ist 0{,}0, d.\,h.\ ein Claim auf \code{inferred} endet als Hypothese.
\end{itemize}

Die Trennung ist die Grundlage für die nachfolgend beschriebene Pipeline -- und zugleich der Punkt, an dem die sechs Werte in der praktischen Wirksamkeit sehr ungleich verteilt sind.

\section{Zweistufiges Claim-Binding}\label{sec:binding}

Die \code{evidence\_binder} und \code{evidence\_entailment} (in \code{backend/app/services/}) implementieren ein zweistufiges Binding:

\begin{enumerate}
  \item \textbf{Retrieval-Stage:} Cosine-Ähnlichkeit zwischen Claim-Vektor und Evidence-Vektor. Schwellen 0{,}55 (niedrige Confidence) und 0{,}65 (hohe Confidence).
  \item \textbf{Entailment-Stage:} Regelbasierte deterministische Prüfung -- SUPPORTED, CONTRADICTED, RELATED\_ONLY oder INSUFFICIENT. Diese Stufe prüft Zahl, Bezugsgruppe, Menge und Modalität explizit vor dem Embedding und verhindert damit das in~\cite{adr-0011} belegte Defektmuster \enquote{61\,\% Zeitersparnis $\to$ 61\,\% bewerten positiv}.
\end{enumerate}

Das \code{supports\_claim}-Feld wird ausschließlich in der Entailment-Stufe gesetzt; dies ist im Code-Kommentar zu \code{evidence\_entailment.py} ausdrücklich festgehalten. Diese Architektur ist zentral für das Verständnis des gesamten Audit-Befundes, weshalb sie in Kapitel~\ref{ch:ergebnisse} erneut aufgegriffen wird.

\section{Anti-Self-Confirmation-Mechanismen}\label{sec:anticonf}

\agora{} implementiert vier strukturelle Anti-Self-Confirmation-Mechanismen, die in der folgenden Tabelle zusammengefasst sind. Ausführlich werden sie in Kapitel~\ref{ch:ergebnisse} bewertet.

\begin{table}[ht]
  \centering
  \small
  \begin{tabular}{@{}p{0.32\textwidth}p{0.58\textwidth}@{}}
    \toprule
    \textbf{Mechanismus} & \textbf{Wirkung} \\
    \midrule
    Echo-Chamber-Index (\code{network\_analytics.py}) &
    Misst intra-cluster-Anteil synthetischer Aktionen; drosselt Confidence bei Werten $> 0{,}75$ via \code{apply\_echo\_cap}. \\
    \midrule
    DACH-Quoten-Kalibrierung &
    Verhindert willkürliche Persona-Population; IT-Anteil hard-gecappt $\leq 12\,\%$; DACH-Branchenanteile gegen Destatis WZ 2008 / Mikrozensus 2024 kalibriert. \\
    \midrule
    Wording-Glossar (\code{\_red\_team\_system\_prompt}) &
    Verbietet Vorhersagesprache; Volltext-Snapshot; Red-Team-Enforcement post-Assembly. \\
    \midrule
    Confidence-Capping &
    Caps für $< 2$ Evidence-Items (0{,}59), \code{model\_generated\_inference}-Gewicht 0{,}0; Echo-Cap auf 0{,}84/medium. \\
    \bottomrule
  \end{tabular}
  \caption{Anti-Self-Confirmation-Mechanismen in \agora{}; Wirkung und Code-Anker.}
  \label{tab:anticonf}
\end{table}

\section{Simulations-Subprozess}\label{sec:simulation}

Der Simulations-Subprozess integriert OASIS- bzw.\ CAMEL-basierte soziale Simulation. Eine Population synthetischer Personas wird über bis zu 5500~LOC Subprozess-Code über mehrere Runden hinweg sozial aktiv (Posts, Comments, Reshares). Der Echo-Chamber-Index berechnet sich als Verhältnis der Intra-Cluster-Aktionen zur Gesamtzahl und koppelt an das Confidence-Gating. Die Simulation ist \emph{code-verifiziert nicht deterministisch} (\code{random.*} ohne \code{random.seed()} an mehreren Stellen; einzige Determinismus-Insel ist \code{louvain(seed=42)} im Community-Detection-Schritt).

\section{Vom Code zum Paper: Annotationsstrategie}\label{sec:annotations}

Jede Code-Aussage in den folgenden Kapiteln ist über einen der vier Schlüssel\texttt{\code{C<n>}}, \texttt{\code{A<n>}}, \texttt{\code{I<n>}} und \texttt{\code{R<n>}} belegt. Eine vollständige Tabelle der Belege ist in der Citation-Registry~\cite{citation-registry} hinterlegt; die Datei listet 57~Code-Belege, 5~ADRs, 12~Issues, 4~Referenzläufe und 14~Literaturschlüssel. Im Text werden die Code-Belege als \texttt{\src{C11}{split\_text ohne Manifest, \code{graph\_build.py:439}}} inline nachgewiesen; \code{AGENTS.md} und \code{README.md} werden ausdrücklich \emph{nicht} als Belege herangezogen.